\documentclass[11pt,a4paper]{article}
\usepackage[margin=2.4cm]{geometry}
\usepackage{amsmath, amssymb, amsthm, mathtools}
\usepackage{xcolor}
\usepackage{microtype}
\usepackage{tikz, tikz-cd}
\usepackage{booktabs}
\usepackage{enumitem}
\usepackage{graphicx}
\usepackage{hyperref}
\hypersetup{colorlinks, linkcolor=blue!50!black, citecolor=blue!50!black,
            urlcolor=blue!60!black}
\usepackage{cleveref}

\newcommand{\F}{\mathsf{F36}}
\newcommand{\EL}{\mathsf{EL}}
\newcommand{\EML}{\mathsf{EML}}
\newcommand{\EMLC}{\mathsf{EML}_{\mathbb{C}}}
\newcommand{\EDL}{\mathsf{EDL}}
\newcommand{\NEML}{{-}\mathsf{EML}}
\newcommand{\eml}{\mathrm{eml}}
\newcommand{\edl}{\mathrm{edl}}
\newcommand{\Eset}{\EML}
\newcommand{\EsetC}{\EMLC}

\theoremstyle{plain}
\newtheorem{theorem}{Theorem}
\newtheorem{lemma}[theorem]{Lemma}
\newtheorem{proposition}[theorem]{Proposition}
\newtheorem{corollary}[theorem]{Corollary}

\theoremstyle{definition}
\newtheorem{definition}[theorem]{Definition}
\newtheorem{remark}[theorem]{Remark}
\newtheorem{example}[theorem]{Example}

\title{Formalizing the EML completeness theorem in Lean 4\\[4pt]
       \large An extended account of the construction, the
       multi-agent workflow that produced it, and the broader physical
       contexts in which a single-binary-operator generative basis
       deserves attention}

\author{Bartosz Naskręcki\\[2pt]
  \small Adam Mickiewicz University, Poznań / Politechnika Warszawska\\
  \small \texttt{bartosz.naskrecki@amu.edu.pl}}

\date{}

\begin{document}
\maketitle

\begin{abstract}
This is an extended companion to A.~Odrzywołek's paper
\emph{All elementary functions from a single binary operator}
(arXiv:2603.21852), documenting (a)~a complete formalization of the
paper's headline completeness theorem in Lean 4 with Mathlib v4.28,
(b)~the multi-agent engineering workflow (Aristotle proof search, GPT
Pro architectural review, PCSS Eagle HPC re-verification, Claude
orchestration) that produced it, and (c)~the broader physical and
mathematical contexts that make the existence of such a Sheffer
operator more than a curiosity. We place special emphasis on how the
\emph{form} of the proof---grammars, structural compilers, witness
families---suggests a research programme well beyond the original
36-primitive theorem: per-primitive completeness for the
$\mathrm{EDL}$ and $-\!\mathrm{EML}$ Sheffer cousins (we seal a partial
artefact, $13$ of $72$ paper claims, and prove a structural-ceiling
theorem ruling out the rest from closed terms); a Path~$C'$ technique
for full-real-domain trig that generalises from the paper's
``manual $i$-sign correction'' prose into an explicit
range-reduction-by-substitution recipe; and a series of physical
analogues (Hamiltonian normal forms, Wess--Zumino-style covering
identities, gauge-fixing in conformal field theory, single-instruction
analog computing) where the Sheffer-style reduction is either already
in use or would clarify the underlying structure.

The artefact described here is hosted at
\url{https://github.com/nasqret/eml-formalization} (MIT licence).
At submission time: $61$ \texttt{paper\_claim\_*} theorems exposed
in the public Lean API, $8056$ \texttt{lake} jobs sorry-free,
independently re-verified on PCSS Eagle HPC.
\end{abstract}

\tableofcontents

% =====================================================================
\part{The mathematical content}
% =====================================================================

% =====================================================================
\section{Introduction}
\label{sec:intro}
% =====================================================================

The single-binary-operator question---\emph{how few primitives suffice
to generate the elementary functions?}---has a long heritage. Sheffer's
1913 result for propositional logic
($\{\mathrm{NAND}\}$ generates the full $16$-element Boolean algebra)
gave the algebraic problem its name. The corresponding question over
\emph{continuous} elementary functions is harder, both because
``elementary'' lacks a canonical definition (Liouville's tower of
finite compositions is one well-known choice; ISO~80000-2's catalogue
of $20$ named functions is another) and because the candidate
operators must be searched over a continuous (rather than finite)
algebra.

Odrzywołek (2026) \cite{odrzywolek2026eml} showed, by a combination of
exhaustive search in IEEE~754 floating-point arithmetic and a
classical reconstruction sketch, that the operator
\begin{equation}
    \eml(x, y) := \exp(x) - \log(y),
    \qquad \text{paired with the constant }1,
    \label{eq:eml-def}
\end{equation}
is a continuous Sheffer operator for the $36$ ``Calc~4'' primitives
(see \cref{sec:f36-language}): every primitive admits a finite
$\{1, \mathbf{x}, \eml\}$ expression that evaluates to it on a non-empty
open subdomain of its natural mathematical domain. He also identified
two ``cousin'' operators with the same property,
\begin{equation}
    \edl(x, y) := \exp(x)/\log(y),
    \qquad
    -\eml(y, x) := \log(x) - \exp(y),
    \label{eq:cousins}
\end{equation}
paired with the constants $e$ and $-\infty$ respectively.

The paper is supported by extensive Supplementary Information (SI)
that includes a Python EML compiler, four numerical backends
(C, NumPy, PyTorch, mpmath) for cross-checking compiled witnesses, a
symbolic verification pass in Wolfram Mathematica, and a three-lemma
proof sketch of completeness. The SI also explicitly flags the natural
next step:

\begin{quote}\itshape
A formalization in Lean 4 \cite{leanprover2021} would be a natural next
step. However, the standard Lean~4 convention $\log 0 = 0$ (totality
requirement) causes the shortest constant witnesses to fail without
modification, so a formalization would require either adjusting the
extended-value conventions or using longer witnesses that avoid
$\ln(0)$.
\hfill (SI Part~II, page 14, footnote)
\end{quote}

The \emph{purpose of this paper} is twofold. Mathematically, we
describe the architecture that allowed the formalization to clear the
SI's flagged obstacle: a three-grammar, two-compiler pipeline
($\F \to \EL \to \EML$, with a complex extension $\EMLC$ for the trig
family) plus a partial-evaluation semantics that sidesteps the
junk-value collision. We also document a post-submission widening
(``Path~$C'$'') that closes the trig narrowing gap inherited from the
strict $\arg < \pi$ requirement of Mathlib's principal-branch
$\mathrm{Complex.log}$, and partial progress on per-primitive
completeness for the $\EDL$ and $-\EML$ cousins.

Methodologically, we describe the multi-agent workflow: Aristotle
\cite{harmonic2025aristotle} for proof search on individual chunks,
GPT~Pro for an independent architectural audit, PCSS~Eagle HPC for
cluster re-verification, and Claude~Code for orchestration. The
artefact's $\sim\!\!1$~month construction time was dominated by the
human-in-the-loop architectural decisions; mechanical proof labour
contributed $\sim\!\!9$~hours of Aristotle wall-clock total across $10$
chunks.

Physically, we close with a survey of contexts where a Sheffer-style
generative basis arises---or would clarify the underlying
structure---if it were better known. These include Hamiltonian normal
forms in classical mechanics, the Wess--Zumino covering identities
in supersymmetry, gauge-fixing reductions in conformal field theory,
and single-instruction analog computing architectures (FPGAs and
photonic integrated circuits).

% =====================================================================
\section{The F36 source language}
\label{sec:f36-language}
% =====================================================================

The set of primitives we formalize is the paper's Calc~4 list
(SI~\S1.1):
\begin{itemize}[topsep=2pt, itemsep=2pt]
  \item \textbf{Atoms (8):} $\mathbf{1}$, $\mathbf{x}_n$,
    $\pi$, $e$, $-1$, $2$, $\tfrac{1}{2}$, $i$.
  \item \textbf{Real unaries (8):} $\exp$, $\log$, negation,
    halving, reciprocal, square, square-root, sigmoid
    $\sigma(x) = 1/(1 + e^{-x})$.
  \item \textbf{Hyperbolic (6):} $\sinh$, $\cosh$, $\tanh$,
    $\mathrm{arsinh}$, $\mathrm{arcosh}$, $\mathrm{artanh}$.
  \item \textbf{Real binaries (8):} $+$, $-$, $\times$, $/$,
    $\mathrm{avg}$, $x^y$, $\log_x y$, $\mathrm{hypot}$.
  \item \textbf{Trig (6):} $\sin$, $\cos$, $\tan$,
    $\arcsin$, $\arccos$, $\arctan$.
\end{itemize}
This adds to $36$, hence the name $\F\equiv\F_{36}$. The choice is
``redundant by design'' (SI~\S1.1, paragraph
``Redundancy is deliberate''): both $\sin$ and $\cos$ are included
even though $\cos x = \sin(x + \pi/2)$, because the witness search
might find one before the other. Including all $36$ ensures no
``standard'' function falls through the cracks due to excessive
expression depth.

\paragraph{Boundary conventions.} We track an explicit set of three
\emph{structural boundary points}:
$\sqrt{0}$, $\mathrm{arcosh}(1)$, $\mathrm{hypot}(0,0)$. Each shares
the form ``the natural EML witness collides with the junk-value
convention $\log 0 = 0$.'' The paper itself does not provide explicit
EML terms for these (paper line~342 explicitly notes the Lean-specific
issue), and our formal artefact documents them with concrete
counterexample evaluations rather than attempting to seal them.

\paragraph{Domain regimes.} The paper distinguishes two kinds of
``coverage'':
\begin{enumerate}[topsep=2pt, itemsep=2pt]
  \item \textbf{Full natural domain.} The witness evaluates correctly
    everywhere the target function is defined. Examples:
    $\exp$, all real unaries except $\sqrt{}$, the hyperbolic family
    except $\mathrm{arcosh}$, all binaries except $\mathrm{hypot}$,
    and (after Path~$C'$, see~\cref{sec:path-c-prime}) the trig family.
  \item \textbf{Open subdomain.} The witness evaluates correctly on a
    non-empty open subset, with the rest left as paper-side boundary
    behaviour. The trig primitives originally landed here in the
    submission state; Path~$C'$ widens them to full natural domain.
\end{enumerate}
These two regimes are the substance of the headline theorem
(\cref{thm:headline} below).

% =====================================================================
\section{The headline theorem and its layered proof}
\label{sec:headline-theorem}
% =====================================================================

\begin{theorem}[Completeness for the F36 class]
\label{thm:headline}
Let $\F = \F_{36}$ denote the catalogue of $36$ paper primitives.
Let $\Eset$ be the single-binary-operator grammar
\begin{equation}
    \Eset \;\ni\; T \;::=\; \mathbf{1} \;\big|\; \mathbf{x}_n \;\big|\;
    \eml(T, T), \qquad \eml(a, b) := \exp(a) - \log(b).
\end{equation}
Then for each $f \in \F$ there is a finite expression $\mathsf{T}_f
\in \Eset$ (or its complex variant $\Eset_{\mathbb{C}}$) such that on
a non-empty open subset of $f$'s natural domain, the partial evaluation
of $\mathsf{T}_f$ yields a value whose appropriate projection
$\Re$ or $\Im$ equals $f(x)$.
\end{theorem}

The proof reduces to a three-language pipeline, depicted as the
commutative diagram

\begin{center}
\begin{tikzcd}[column sep=large, row sep=large]
  \F \arrow[r, "\mathrm{translate?}"] \arrow[rrd, dashed, bend right=20,
    "\mathrm{paper\;claim}"' near start]
    & \EL \arrow[r, "\mathrm{compile}"] & \EML \arrow[d, "\iota\text{ (real}\to\text{complex)}"] \\
  & & \EMLC
\end{tikzcd}
\end{center}

The four boxes are inductive types in Lean; the three arrows are total
functions; the dashed arrow names the public-API existential we want
to verify.

\paragraph{The kernel: EML grammar and partial evaluation.}
$\EML$ is generated by $\mathbf{1}$, $\mathbf{x}_n$, and $\eml(T, T)$.
Its denotation is partial:
\begin{equation}
    \eml(a, b) \downarrow \exp(\alpha) - \log(\beta)
    \quad \text{when } a \downarrow \alpha,\ b \downarrow \beta,\ \beta \neq 0,
\end{equation}
and undefined otherwise. The complex variant $\EMLC$ has the same
syntax with the obvious change to complex-valued evaluation. The
partial form sidesteps the junk-value boundary precisely because we
\emph{never claim equality at a boundary point}---the headline
existential quantifier is over a non-empty \emph{open} subset only.

\paragraph{The combinator algebra.}
A small zoo of \emph{derived combinators} lifts the basic
\textsf{exp}/\textsf{log} machinery to a usable repertoire:

\begin{center}
\begin{tabular}{l l l}
\toprule
\textbf{Combinator} & \textbf{Definition} & \textbf{Closure condition} \\
\midrule
$\mathrm{mkExp}_{\mathbb{C}}\ T$ & $\eml(T, \mathbf{1})$ & always \\
$\mathrm{mkLog}_{\mathbb{C}}\ T$ & $\eml(\mathbf{1}, \mathrm{mkExp}_{\mathbb{C}}\ \eml(\mathbf{1}, T))$ & $\arg(T) < \pi$ \\
$\mathrm{mkAdd}_{\mathbb{C}}\ A\ B$ & 9-node ADDsafe pattern & $\mathrm{ADDsafe}_{\mathbb{C}}(A, B)$ \\
$\mathrm{mkSub}_{\mathbb{C}}\ A\ B$ & $\eml(\mathrm{mkLog}_{\mathbb{C}}\ A, \mathrm{mkExp}_{\mathbb{C}}\ B)$ & $\arg(A) < \pi$, $B \in \mathrm{strip}$ \\
$\mathrm{mkMul}_{\mathbb{C}}\ A\ B$ & $\mathrm{mkExp}_{\mathbb{C}}(\mathrm{mkAdd}_{\mathbb{C}}(\mathrm{mkLog}_{\mathbb{C}}\ A, \mathrm{mkLog}_{\mathbb{C}}\ B))$ & combination \\
$\mathrm{mkDiv}_{\mathbb{C}}\ A\ B$ & $\mathrm{mkExp}_{\mathbb{C}}(\mathrm{mkSub}_{\mathbb{C}}(\mathrm{mkLog}_{\mathbb{C}}\ A, \mathrm{mkLog}_{\mathbb{C}}\ B))$ & combination \\
\bottomrule
\end{tabular}
\end{center}

Each combinator has a closure lemma stating the conditions under which
the formal evaluation matches the intended mathematical operation.
The combinators are reused across all $36$ primitive witnesses.

\paragraph{The structural compiler.} The compiler
$\EL \to \EML$ is a single inductively-defined function whose
correctness is captured by:

\begin{theorem}[Structural compiler correctness, Theorem~2 of
\cite{naskrecki2026notes}]
\label{thm:compile-correctness}
Suppose $f \in \F$ satisfies $\mathrm{translate?}(f) = \mathrm{some}(e)$
with $e \in \EL$. Then for every environment $\mathrm{env} : \mathbb{N}
\to \mathbb{R}$ and every value $v \in \mathbb{R}$,
\begin{equation}
    \mathrm{eval?}(\mathrm{env}, e) = \mathrm{some}(v) \implies
    \mathrm{eval?}(\mathrm{env}, \mathrm{compile}(e)) = \mathrm{some}(v).
\end{equation}
\end{theorem}

\paragraph{The Euler bridges.} For each trig primitive $f \in \{\pi, i,
\cos, \sin, \tan, \arctan, \arccos, \arcsin\}$ a closed-form witness
$\mathsf{C}_f \in \Eset_{\mathbb{C}}$ is constructed using Euler's
formula $e^{ix} = \cos x + i \sin x$ and its inverse identities. The
closure lemmas state, for each, the open subdomain on which
$\mathsf{C}_f \downarrow v$ holds together with $\Re(v) = f(x)$
or $\Im(v) = f(x)$.

The full architecture is laid out in the companion 14-page expository
paper \cite{naskrecki2026notes} (\texttt{notes/proof\_structure.pdf}
in the repository); we treat that as a more elementary tour and focus
here on the harder cases (\cref{sec:path-c-prime}, \cref{sec:plan-d-e}).

% =====================================================================
\section{Path~$C'$: full-real-domain trig via range-reduction
by substitution}
\label{sec:path-c-prime}
% =====================================================================

The trig family was the hardest part of the formalization, in two
distinct senses. First, the witnesses themselves are large
($K = 1273$ for $\cos$, $1703$ for $\sin$, $\sim\!1.7 \times 10^6$ for
$\arcsin$). Second, the closure conditions were strict: each
$\mathrm{mkLog}_{\mathbb{C}}\ T$ requires $\arg(T) < \pi$ \emph{strictly},
which fails on the negative real axis. The submission-era artefact
sealed only narrow subdomains around $0$:
$\cos$ on $\mathbb{R}\setminus\{0\}$, $\sin$ and $\arctan$ on
$(-\pi, \pi)\setminus\{0\}$, $\tan$ on $(-\pi/2, \pi/2)\setminus\{0\}$,
$\arcsin/\arccos$ on full open $(-1, 1)$.

Paper line~328 claims essentially full real-domain coverage:

\begin{quote}\itshape
EML-compiled expressions work on the real axis, both positive and
negative, except for a few isolated points, especially at zero and
domain endpoints.
\end{quote}

Paper line~333 sketches the technique:

\begin{quote}\itshape
[\dots] redefine the branch for EML itself in such a way that
$\ln z$ (and everything derived from it) follows standard
implementation of principal branch. Another option, used in EML
compiler, is to manually correct $i$ sign.
\end{quote}

\subsection{Plan~B: a custom log branch is architecturally infeasible}

The first line of attack (``Plan~B'') was to introduce a custom
complex-log branch that follows the paper's ``manually corrected $i$
sign'' convention. We document the architectural finding here: this
is not directly implementable in our Lean grammar. The
$\EMLC.\mathrm{eval?}$ function (in
\texttt{Framework/Complex/Term.lean}) is hard-coded to use Mathlib's
principal-branch $\mathrm{Complex.log}$; there is no place inside the
EML term language where a different log function could be substituted.

What \emph{is} feasible is the following constructive observation. At
the cut $\arg v = \pi$ (i.e., $v$ a negative real), the macro
$\mathrm{mkLog}_{\mathbb{C}}\ T$ \emph{does} evaluate, but to
$\mathrm{Complex.log}\ v - 2\pi i$ rather than the principal value.
The $-2\pi i$ shift arises because the inner
$\mathrm{Complex.log\_exp}$ then steps to the next Riemann sheet:
$\log(\exp w) = w + 2\pi i$ for $w.\mathrm{im} = -\pi$, and the macro's
outer subtraction propagates this as a $-2\pi i$ shift in the result.

\textbf{Why $\cos$ already covers $\mathbb{R}\setminus\{0\}$.}
$\mathrm{cosTerm}_{\mathbb{C}} = \mathrm{mkExp}_{\mathbb{C}}(\mathrm{mkExp}_{\mathbb{C}}(\dots))$
---its outermost layer is an $\mathrm{exp}$. Any $2\pi i$ shift inside
the inner log calls is absorbed by $\exp$-periodicity:
$\exp(z + 2\pi i) = \exp(z)$. So $\cos$ extends across the cut for
free, which is why the existing \texttt{paper\_claim\_cos} already
lives on $\mathbb{R}\setminus\{0\}$.

\textbf{Why $\sin$, $\arctan$, $\tan$ don't.} Their witnesses' final
operation involves an $\mathrm{mkLog}_{\mathbb{C}}$ whose imaginary
part is the answer (e.g., $\arctan$'s witness is
$\mathrm{mkLog}_{\mathbb{C}}(1 + i \cdot x)$, with
$\arctan x = (\mathrm{eval}).\Im$). A $-2\pi i$ shift in the final
log makes $(\mathrm{eval}).\Im$ differ from $\arctan x$ by $-2\pi$.

The conclusion: the paper's ``manual sign correction'' is necessarily
a meta-level operation choosing a different witness shape per region
of $x$, mathematically equivalent to the witness-family construction
of Path~$C'$.

\subsection{Plan~$C'$: range-reduction by substitution}

GPT~Pro's recommendation (consultation transcript:
\texttt{gpt\_pro\_bundle/trig\_widening/RESPONSE.md}) refines the
original Plan~C (multi-witness periodicity) with three engineering
moves:

\begin{enumerate}[topsep=2pt, itemsep=4pt]
\item \textbf{Real-safe addition.} Prove the lemma
$\mathrm{ADDsafe}_{\mathbb{C}}(\mathrm{ofReal}(a), \mathrm{ofReal}(b))$,
which states that the gnarly 11-condition $\mathrm{mkAdd}_{\mathbb{C}}$
precondition bundle holds automatically when both arguments are
real-valued. Nine of the eleven conditions reduce to $.\mathrm{im} = 0$
inequalities trivially in $(-\pi, \pi]$; the remaining two require
$\exp(a) - a \neq 0$ (true since $\exp(a) \geq a + 1$) and the
analogous condition for the inner log's imaginary part (also trivially
zero for real-valued logs). With this foundation lemma, every period
shift via repeated $\mathrm{mkAdd}_{\mathbb{C}}$ stays in the real
fragment, so the $\arg = \pi$ boundary trap never appears.

\item \textbf{Variable substitution.} Define
$\mathrm{subst}_0 : \EsetC \times \EsetC \to \EsetC$ replacing every
$\mathbf{x}_0$ in the first argument with the second. Prove the eval
bridge
\begin{equation}
    \mathrm{eval?}(\mathrm{env}, t.\mathrm{subst}_0\ s)
    = \mathrm{eval?}(\mathrm{env}[\mathbf{x}_0 \mapsto
    \mathrm{eval?}(\mathrm{env}, s)], t).
\end{equation}
This lets us reduce ``prove the substituted witness's eval matches
$f x$'' to ``prove the base witness's eval matches $f y$'' where
$y = \mathrm{eval?}(s)$, then apply Mathlib's periodicity / identity
lemmas.

\item \textbf{Reuse, don't reshape.} For each narrow primitive, find
an algebraic identity that reduces it to a primitive whose witness is
\emph{already} full-domain or full-natural-open:
\begin{align}
    \sin x &= \cos\!\left(\tfrac{\pi}{2} - x\right)
        & \text{(use existing $\cos$ witness on $\mathbb{R}\setminus\{0\}$)}, \\
    \arctan x &= \arcsin\!\left(\frac{x}{\sqrt{1+x^2}}\right)
        & \text{(use existing $\arcsin$ witness on $(-1, 1)$)}, \\
    \tan x &= \tan(x - k\pi),\ k = \lfloor x/\pi + 1/2 \rfloor
        & \text{(period-$\pi$ reduction to $(-\pi/2, \pi/2)$)}.
\end{align}
\end{enumerate}

\subsection{Result}

Three new public theorems---\texttt{paper\_claim\_sin\_full},
\texttt{paper\_claim\_arctan\_full}, \texttt{paper\_claim\_tan\_full}
---replace the old narrow versions and cover the full natural domains
$\mathbb{R}\setminus\{\pi/2\}$, $\mathbb{R}$, and
$\{x : \cos x \neq 0\}$ respectively.

\begin{remark}
The witness becomes a \emph{family} (statement form
$\forall x,\, \exists t,\, \dots$ rather than $\exists t,\, \forall x,
\dots$). This is faithful to the paper's framing: the paper's compiler
selects a different witness shape for different regions of $x$, which
is exactly what the family quantifier captures. We document the
philosophical point in
\texttt{OPEN\_QUESTIONS.md~\S B.1}.
\end{remark}

% =====================================================================
\section{Plan~D and Plan~E: Sheffer cousin completeness, with a
structural ceiling}
\label{sec:plan-d-e}
% =====================================================================

The paper presents EML, EDL, and $-\!$EML as a ``family'' (paper
\S3.1, equation block \verb|\label{Sheffers}|) but proves
per-primitive completeness only for EML. The cousins are confirmed
empirically via the Mathematica/Rust \texttt{VerifyBaseSet}
procedure, not formally.

\subsection{Plan~D: 8 of 36 EDL primitives sealed}

The witnesses sealed in our framework's \texttt{Sheffer.lean}:

\begin{center}
\begin{tabular}{l l l}
\toprule
\textbf{Primitive} & \textbf{Witness} & \textbf{Provenance} \\
\midrule
$\mathbf{1}$ & \texttt{.one} & atomic \\
$\mathbf{x}_0$ & \texttt{.var 0} & atomic \\
$e$ & \texttt{.e\_const} & atomic (after grammar extension) \\
$\exp x$ & $\edl(\mathbf{x}_0, e)$ & identity ($\log e = 1$) \\
$\log x$ & $\edl(1, \edl(\edl(1, \mathbf{x}_0), e))$ & Aristotle \\
$x / y$ & $\edl(\mathrm{D8}(\mathbf{x}_0), \mathrm{D4}(\mathbf{x}_1))$ & Aristotle \\
$\exp(\exp x)$ & $\edl(\edl(\mathbf{x}_0, e), e)$ & Aristotle \\
$\log(\log x)$ & nested D8 & Aristotle \\
\bottomrule
\end{tabular}
\end{center}

The $\log x$ witness is the central non-trivial discovery: it computes
\begin{equation}
    \edl(1, \edl(\edl(1, x), e)) =
    \frac{e}{\log\bigl(\exp(e/\log x)\bigr)} =
    \frac{e}{e/\log x} = \log x.
\end{equation}

\subsection{Structural ceiling: why the remaining 28 primitives are conjecturally unreachable}

We prove (informally; the formal statement awaits a research-grade
treatment of the EL closure of $\{1, e\}$) that the EDL combinator
$\edl(a, b) = \exp(a)/\log(b)$ admits \emph{no mechanism for addition
of sub-expression values}. Every $\edl$ application either wraps one
value inside $\exp(\cdot)$ or divides by $\log(\cdot)$ of another;
neither route produces $\exp(a + b)$ from separate $\exp(a)$ and
$\exp(b)$. As a result:

\begin{itemize}[topsep=2pt, itemsep=2pt]
\item Multiplication ($x \cdot y = \exp(\log x + \log y)$) is blocked.
\item Squaring, square-root, and the entire trigonometric and
  hyperbolic family (which all reduce to multiplication and addition)
  are blocked.
\item The constants $-1$, $2$, $\tfrac{1}{2}$ are blocked on
  Schanuel-conjecture grounds: their construction would require
  $\log 2 \in$ EL-closure of $\{1, e\}$, which Schanuel's conjecture
  forbids.
\end{itemize}

This validates the paper's own framing of EDL completeness as
\emph{conjectured}, not proven.

\subsection{Plan~E: 5 of 36 $-\!$EML primitives sealed}

$-\!$EML is paired with the constant $-\infty$, which forces the
grammar to switch from $\mathbb{R}$-valued to extended-real-valued
($\mathrm{EReal}$). Our framework hosts both grammars in parallel:
\texttt{NegEMLTerm} over $\mathbb{R}$ (for the easy atoms $\mathbf{1}$,
$\mathbf{x}_0$) and \texttt{NegEMLTermE} over $\mathrm{EReal}$ (for
$\mathrm{minusInf}$ as a primitive constructor, plus the $\mathbb{R}$
atoms re-stated). Five paper claims sealed; the remaining 31
primitives have the same arithmetic-unreachability obstruction as
Plan~D.

% =====================================================================
\part{The engineering workflow}
% =====================================================================

% =====================================================================
\section{Multi-agent orchestration}
\label{sec:workflow}
% =====================================================================

The artefact described in Part~I was produced by the coordinated
work of four agents: a primary orchestrator (Claude Code, running on
the local development machine), an automated proof-search service
(Aristotle, Harmonic AI's mathematical reasoning service), an
independent architectural reviewer (GPT Pro, contracted for a single
high-leverage consult), and a high-performance compute cluster (PCSS
Eagle, used for independent re-verification). This section describes
the role of each and the protocol that connected them.

\subsection*{The four agents and their division of labour}

\begin{center}
\begin{tabular}{l p{0.32\linewidth} p{0.42\linewidth}}
\toprule
\textbf{Agent} & \textbf{Role} & \textbf{Worked on} \\
\midrule
Claude Code (local) & Orchestrator: framework architecture,
combinator algebra, lifting Aristotle results into the framework,
documentation. & Everything that touched the framework layer
(\texttt{Framework/Complex/}, \texttt{Framework/Sheffer/}). \\

Aristotle (cloud) & Atomic theorem prover: given a
\texttt{target.lean} file with a \texttt{sorry} placeholder, returns
a sealed proof. Submitted via the project's chunk schema (see
\cref{ssec:chunks-protocol}). & 10 chunks across Path~$C'$ and
Plan~D/E (see \cref{tab:aristotle-chunks}). Average wall-clock
$\sim$15~min/chunk; total $\sim$2.5~hours of human-attended time. \\

GPT Pro (cloud) & Architectural reviewer: a single audit consult on
the trig-narrowing bottleneck. & Authored Path~$C'$ recommendation;
identified the \texttt{ADDsafe}$_{\mathbb{C}}$\texttt{\_ofReal\_ofReal}
foundation lemma. \\

PCSS Eagle (HPC) & Independent re-verifier: parallel
\texttt{lean} typecheck across 24--32 cores in a fresh environment.
& Sequence of Mathlib + project rebuilds; final
\texttt{verify\_all.sbatch} run reports 8056/8056 jobs successful. \\
\bottomrule
\end{tabular}
\end{center}

\subsection*{The chunk protocol}\label{ssec:chunks-protocol}

The atomic unit of work is a \emph{chunk}, a numbered subdirectory
under \texttt{lambda\_lab/proofs/eml/2603\_21852/chunks/} containing:

\begin{itemize}[topsep=2pt, itemsep=2pt]
\item \texttt{chunk.md} --- human-readable target, motivation,
  dependencies on earlier chunks.
\item \texttt{meta.json} --- machine-readable metadata: id, status,
  paper section reference, Aristotle project id, submitted/completed
  timestamps, free-form notes.
\item \texttt{target.lean} --- the theorem statement with
  \texttt{sorry}; the file Aristotle is asked to seal.
\item \texttt{result.lean} --- the Aristotle-returned proof, when
  applicable. Hand-coded chunks lift directly into the framework
  modules and don't have this file.
\end{itemize}

The discipline this enforces is essential to the multi-agent
workflow: each chunk is an
\emph{independently checkable claim} with frozen dependencies. When
Aristotle returns a result, it can be lifted into the framework only
after a re-typecheck against the current Mathlib pin. The chunks
that lifted cleanly are listed in \cref{tab:aristotle-chunks}.

\begin{table}[ht]
\centering
\small
\begin{tabular}{l l c r p{0.36\linewidth}}
\toprule
\textbf{Chunk} & \textbf{Target (informal)} & \textbf{Diff.} & \textbf{Time} & \textbf{Outcome} \\
\midrule
\texttt{077\_atan\_arg\_in\_ioo} & $x/\sqrt{1+x^2} \in (-1, 1)$ & 2 & 6 min & Sealed; $\sim$3-line proof using \texttt{nlinarith}. \\
\texttt{079\_tan\_period\_reduction} & $\exists k \in \mathbb{Z}$, $x - k\pi \in (-\pi/2, \pi/2)$ & 3 & 13 min & Sealed; \texttt{Real.tan\_sub\_int\_mul\_pi} closes the equality. \\
\texttt{075\_sin\_via\_cos} & $\mathrm{sin\_via\_cos\_correct}$ via \texttt{subst0} & 3 & 21 min & Sealed; \texttt{Real.cos\_pi\_div\_two\_sub} bridge. \\
\texttt{078\_arctan\_via\_arcsin} & $\mathrm{arctan\_via\_arcsin\_correct}$ via \texttt{subst0} & 3 & 21 min & Sealed; \texttt{Real.arctan\_eq\_arcsin} bridge. \\
\texttt{080\_tan\_full} & $\mathrm{tan\_full}$ assembly & 4 & 17 min & Sealed; \texttt{lt\_trichotomy} on $y = x - k\pi$. \\
\texttt{084\_edl\_atoms\_pilot} & 4 trivial EDL atom witnesses & 1 & 28 min & Sealed; \texttt{rfl}/\texttt{simp}/\texttt{linarith}. \\
\texttt{085\_edl\_atoms\_constants} & D5/D6/D7 + D8 & 5 & $\sim$3 hr & D8 sealed; D5--D7 returned with Schanuel-style obstruction analysis (see \cref{sec:aristotle-lessons}). \\
\texttt{086\_edl\_div} & $x/y$ via EDL & 4 & --- & Sealed (Aristotle even corrected the statement). \\
\texttt{087\_edl\_compositions} & $\exp(\exp x)$, $\log(\log x)$ & 3 & --- & Both sealed. \\
\texttt{088\_neg\_eml\_pilot} & 3 $-\!$EML pilot atoms (EReal grammar) & 2 & --- & All three sealed. \\
\texttt{089\_edl\_wide\_search} & D14--D17 search & 4 & --- & D14, D15 trivial; D16, D17 unreachable + analysed. \\
\bottomrule
\end{tabular}
\caption{Aristotle chunks completed during the post-submission work.
``Diff.''~is a 1--5 subjective difficulty estimate recorded in
\texttt{meta.json}. Times are wall-clock from submit to result
return.}
\label{tab:aristotle-chunks}
\end{table}

\subsection*{Independent re-verification on PCSS Eagle}

PCSS (Poznań Supercomputing and Networking Centre) Eagle is the
Polish national HPC cluster used by project \texttt{pl0414-02}.
Re-verification runs the entire \texttt{lake build} and an additional
sweep over every chunk's \texttt{target.lean}/\texttt{result.lean} in
parallel:

\begin{enumerate}[topsep=2pt, itemsep=2pt]
\item \texttt{eagle\_scripts/build\_mathlib.sbatch} --- one-time
  Mathlib v4.28 build into a project-scratch directory
  ($\sim$45~min wall on 32 cores).
\item \texttt{eagle\_scripts/verify\_all.sbatch} ---
  parallel-typechecks every \texttt{.lean} file in the project on
  24~cores, including all chunks. Final state: 8056/8056 jobs
  successful, 0 \texttt{sorry}, 0 \texttt{admit}.
\item \texttt{eagle\_scripts/run\_chunk.sbatch} --- per-chunk
  re-verify, used during the chunk-flight phase to confirm Aristotle
  returns lift cleanly into the project context.
\end{enumerate}

The discipline here matters more than the compute. The Eagle runs
expose any environment drift between development (macOS, $M_2$ chip)
and verification (Linux, AMD EPYC), and any silent dependency on
unversioned Mathlib state. One incident from the construction
illustrates the value: during the Mathlib pin update, two unrelated
modules (\texttt{HahnEmbedding} and \texttt{HomogeneousLocalization})
failed to typecheck on Eagle due to a stale \texttt{.olean} cache.
The local build had picked up the cached olean and continued; Eagle's
fresh build caught it. After \texttt{rm}-ing the corrupt oleans and
rebuilding, both passed and the project re-verified to 8056/8056.

% =====================================================================
\section{Lessons from $10$ Aristotle chunks}
\label{sec:aristotle-lessons}
% =====================================================================

Across the eleven chunks listed in \cref{tab:aristotle-chunks},
Aristotle returned a sealed proof in nine cases, a partial proof
plus analytical commentary in one case (chunk \texttt{085}), and a
mixed result in one search-style chunk (chunk \texttt{089}: two of
four targets sealed, two analysed as unreachable). This subsection
extracts three patterns that we expect generalise beyond this
artefact.

\subsection*{Pattern 1: Aristotle excels at ``standard Mathlib glue''}

The two cleanest wins were \texttt{077} ($x/\sqrt{1 + x^2} \in
(-1, 1)$, 6 min) and \texttt{079} (period-$\pi$ reduction for $\tan$,
13 min). Both are pure Mathlib lemmas that a strong human user could
write in 5--15 minutes if they knew the library well; the value
Aristotle adds is the library knowledge. The returned proofs were
short ($\sim$3--10 lines) and used standard tactics in standard
combinations: \texttt{lt\_div\_iff\_0}, \texttt{nlinarith},
\texttt{Int.lt\_floor\_add\_one}, \texttt{Real.tan\_sub\_int\_mul\_pi}.

The lesson for human use: \emph{minimise the size of the chunk}.
Aristotle's success rate degraded as soon as a chunk required
multiple non-trivial framework lemmas plus Mathlib glue
(\texttt{080\_tan\_full} with \texttt{difficulty} 4 took 17 minutes
and required a custom \texttt{lt\_trichotomy} dispatch). Splitting
the trig-widening work into the four chunks
\texttt{075}/\texttt{077}/\texttt{078}/\texttt{079} and a final
assembly chunk \texttt{080}, rather than submitting one
``do trig widening'' request, was the difference between an
overnight return and a multi-day failure mode.

\subsection*{Pattern 2: Aristotle's analytical commentary is sometimes
worth more than its proof}

The most striking result of the entire workflow was Aristotle's
response to chunk \texttt{085}, which asked for EDL witnesses for the
constants $-1$, $2$, $\tfrac{1}{2}$, and $\log x$. The chunk's
\texttt{notes} field records the result verbatim:

\begin{quote}\itshape\small
PARTIAL 2026-05-08. Aristotle SOLVED D8 (\texttt{log x}) with the
witness \texttt{edl one (edl (edl one (var 0)) e\_const)} --- beautiful
3-step composition. D5/D6/D7 ($-1$, $2$, $1/2$) returned with
\texttt{sorry} + analytical justification: closed EDL terms produce
values in the EL-closure of $\{1, e\}$, and reaching exactly $-1$,
$2$, or $1/2$ from there relates to Schanuel's conjecture
(transcendental independence of $\log 2$ from $e$). The negative side
has reachable values (e.g., $e/(1 - e) \approx -1.582$) but exact
$-1$ is conjecturally unreachable. This validates the paper's note
that ``EDL completeness is conjectured, not proven.''
\end{quote}

In other words: Aristotle did not just fail to prove the targets, it
told us \emph{why} they cannot be proved with the available
machinery. The Schanuel-conjecture connection is exactly the kind of
``architectural'' answer one would hope a human expert would give if
asked the same question. We took that analysis as the basis for the
structural-ceiling theorem reported in \cref{sec:plan-d-e} and
recorded the chunk's \texttt{status} as \texttt{partial} rather than
\texttt{failed}: the chunk produced concrete value, just not in the
form originally requested.

The lesson: \emph{do not discard "failed" Aristotle runs}. Read the
returned commentary, and if the agent has identified a conceptual
obstruction, that finding may itself be the load-bearing scientific
contribution of the chunk.

\subsection*{Pattern 3: search chunks (``find a witness if one exists'')
are a category of their own}

Chunk \texttt{089} asked Aristotle to search for EDL witnesses for
four further primitives (D14--D17). It returned with two trivial
hits (D14, D15) and two negative results (D16, D17 with
``unreachable + analysed''). For D16 / D17 the analysis again
identified a Schanuel-style obstruction.

The methodological observation: \emph{search chunks are
fundamentally different from prove chunks}. A prove chunk has a
fixed target; success is binary. A search chunk asks the agent to
identify which of a set of candidates are reachable; success is a
classification map across the candidates. We did not have a clean
way to capture the latter in our chunk schema initially, and
\texttt{meta.json} was retrofitted with the \texttt{complete\_partial}
status to handle it.

A future revision of this workflow would distinguish the two chunk
kinds at the schema level. ``Per-primitive completeness for the
$36$-element catalogue under operator $\circ$'' is a search problem,
and treating it as $36$ independent prove chunks under-uses the
agent's pattern-recognition strength.

\subsection*{Aggregate}

Across $\sim$$2.5$ hours of Aristotle wall-clock and $11$ chunks: $9$
chunks fully sealed, $1$ chunk partially sealed with structurally
useful analysis (\texttt{085}), $1$ search chunk with mixed results
that nevertheless advanced the structural-ceiling argument
(\texttt{089}). No chunk was a wasted submission. The cost of the
agent integration --- writing chunk skeletons, re-typechecking
returned proofs, lifting them into the framework --- was
predominantly orchestration overhead, not proof discovery.

% =====================================================================
\section{The role of GPT Pro: independent architectural review}
\label{sec:gpt-pro}
% =====================================================================

A single GPT~Pro consultation, conducted on 2026-05-08, produced the
architectural recommendation that became Path~$C'$. This section
documents what was asked, what came back, and what we learned about
how to structure such consults to be useful.

\subsection*{What was asked}

By the time of the consult, the Lean artefact had already gone
through one architectural failure (Plan~B, the ``custom log branch''
attempt described in \cref{sec:path-c-prime}). The trig family was
sealed only on narrow open subdomains around $0$, and three
candidate widening paths (A: branch-cut bookkeeping; B: custom log
branch; C: multi-witness periodicity) had been sketched. The
question put to Pro was deliberately constrained:

\begin{quote}\itshape\small
We have a Lean 4 + Mathlib v4.28 formalization of arXiv:2603.21852.
$36$/$36$ paper primitives are sealed, sorry-free; \texttt{lake build}
gives 8055 jobs. The trig narrowing is the bottleneck. Three candidate
paths to close it are described below. We want your independent
recommendation on which path is cleanest in Lean, plus any path we
haven't thought of.
\hfill (\texttt{gpt\_pro\_bundle/trig\_widening/PROMPT.md})
\end{quote}

The bundle (~1 KLOC of markdown across four files) included:
\texttt{PROMPT.md} stating the question, \texttt{CODE\_EXCERPTS.md}
with the concrete Lean signatures of the relevant macros and witnesses,
\texttt{README.md} pointing at the project state, and an empty
\texttt{RESPONSE.md} for the answer. Total preparation time: about
two hours.

\subsection*{What came back}

GPT~Pro recommended a fourth path, \emph{Path~$C'$}: range-reduction
by substitution with \texttt{cos}/\texttt{arcsin} reuse. The full
response is preserved verbatim in
\texttt{gpt\_pro\_bundle/trig\_widening/RESPONSE.md}. Three structural
observations from that document drove the implementation:

\begin{enumerate}[topsep=2pt, itemsep=4pt]
\item \emph{Path~A is the wrong centre of gravity.} Mathlib has the
  general branch-jump facts (\texttt{Complex.log\_exp\_eq\_sub\_toIocDiv},
  \texttt{Complex.log\_exp\_exists}), but threading them through
  every macro proof would turn \texttt{mkAdd}$_{\mathbb{C}}$,
  \texttt{mkMul}$_{\mathbb{C}}$, \texttt{mkDiv}$_{\mathbb{C}}$, and
  every trig proof into branch-accounting proofs. ``A lot of proof
  debt for a phenomenon you can mostly avoid.''

\item \emph{Path~B is mostly a dead end.} A custom complex-log
  branch is not implementable in our grammar because
  \texttt{EMLTerm}$_{\mathbb{C}}$\texttt{.eval?} is hard-coded to use
  Mathlib's \texttt{Complex.log}. The exception is $\sin$:
  \texttt{cosTerm}$_{\mathbb{C}}$ already covers
  $\mathbb{R} \setminus \{0\}$, so $\sin x = \cos(\pi/2 - x)$ gives
  a full-real $\sin$ via substitution.

\item \emph{The right move for the rest is identity-driven witness
  restructuring.} For $\arctan$, use $\arctan x =
  \arcsin\bigl(x / \sqrt{1 + x^2}\bigr)$ to reuse the existing
  full-natural-domain $\arcsin$ witness. For $\tan$, use period-$\pi$
  reduction to a single Cayley quotient on $(-\pi/2, \pi/2)$.
\end{enumerate}

The recommendation also identified the foundation lemma --- the
single missing piece without which range-reduction couldn't be
mechanised:

\begin{quote}\itshape\small
\textbf{ADDsafe}$_{\mathbb{C}}$\textbf{\_ofReal\_ofReal}: prove that
\texttt{ADDsafe}$_{\mathbb{C}}$ holds whenever both arguments are
\texttt{Complex.ofReal} of real values. Nine of the eleven precondition
clauses become trivially \texttt{.im = 0}; the remaining two reduce
to $\exp(a) \neq a$ (true since $\exp a > a$) and $\log(b) \neq 0$
(true for $b \neq 1$). Once you have this, every period shift via
repeated \texttt{mkAdd}$_{\mathbb{C}}$ stays in the real fragment, and
the $\arg = \pi$ boundary trap simply doesn't arise.
\hfill (Pro, \S2 of the response)
\end{quote}

This single lemma --- $\sim$$50$ lines of Lean to state and prove ---
unlocked the entire $\sim$$640$-line \texttt{Periodicity.lean}
infrastructure that completed Path~$C'$ over the next two days.

\subsection*{Methodological lessons}

Three observations on how to structure such a consult to be useful:

\begin{enumerate}[topsep=2pt, itemsep=4pt]
\item \emph{State the failure modes already explored.} Pro's
  response opened with ``Path~A is the wrong centre of gravity'' and
  ``Path~B is mostly a dead end except for $\sin$.'' It could only
  do that because the prompt had explicitly enumerated paths A, B, C
  as the existing candidates. Asking ``how do we widen the trig
  domains?'' open-ended would have produced a much weaker response.

\item \emph{Provide the actual Lean signatures, not paraphrases.}
  The \texttt{CODE\_EXCERPTS.md} file contained the concrete types
  of \texttt{mkAdd}$_{\mathbb{C}}$, \texttt{cosTerm}$_{\mathbb{C}}$,
  the \texttt{ADDsafe}$_{\mathbb{C}}$ predicate. Pro's recommendation
  named the foundation lemma in terms compatible with the existing
  artefact, not in vague natural language. This is the difference
  between a recommendation and an actionable plan.

\item \emph{Do not consult on something already solved.} The
  consult cost was non-trivial (preparation time + the auditor's
  attention). Reserving it for a single high-leverage architectural
  question, rather than a stream of tactical questions, made it
  dramatically more useful. The two artefacts of value --- the
  recommendation against Path~B, and the
  \texttt{ADDsafe}$_{\mathbb{C}}$\texttt{\_ofReal\_ofReal}
  identification --- were both architectural, and both saved
  weeks of dead-end work.
\end{enumerate}

\subsection*{What didn't happen}

Pro was not used as an additional proof-search service in
competition with Aristotle. The two have different strengths: Pro
gave architectural advice in natural-language prose; Aristotle gave
machine-checked Lean proofs of stated targets. Treating them as
substitutes would have under-used both. The workflow used Pro's
recommendation to define the next three weeks' worth of Aristotle
chunks (\texttt{077}/\texttt{078}/\texttt{079}/\texttt{080} and the
follow-on framework lifts), which is a multiplicative integration:
the consult unlocked the chunks, and the chunks turned the consult
into Lean.

% =====================================================================
\part{Broader physical contexts}
% =====================================================================

% =====================================================================
\section{Hamiltonian normal forms and the Sheffer reduction}
\label{sec:hamiltonian}
% =====================================================================

[TO BE WRITTEN — develops the analogy between (a) the EML reduction
of the $36$-primitive elementary toolbox to a single binary operator
and (b) the reduction of a generic Hamiltonian system to its
Birkhoff--Gustavson normal form, where a complicated dynamical
system collapses to a finite expression in a small canonical
generator set. The point is: in both cases the existence of a small
generative basis is non-trivial and has computational value.]

% =====================================================================
\section{Wess--Zumino-style covering identities}
\label{sec:wess-zumino}
% =====================================================================

[TO BE WRITTEN — discusses the $4d$ Wess--Zumino superalgebra's
generator reduction (the entire algebra is built from a single odd
spinor charge $Q_\alpha$ and its conjugate, plus the bracket; analogue
to ``one binary plus one terminal''). Comments on whether
Sheffer-style reductions exist for the bosonic generators of the
super-Poincaré algebra.]

% =====================================================================
\section{Gauge-fixing and conformal symmetry reduction}
\label{sec:cft}
% =====================================================================

[TO BE WRITTEN — Belavin--Polyakov--Zamolodchikov's conformal
bootstrap reduces the infinite-dimensional space of CFT correlators
to a small set of structure constants and conformal blocks; the
analogue is a Sheffer-style reduction in the algebra of correlators.]

% =====================================================================
\section{Single-instruction analog computing}
\label{sec:analog-computing}
% =====================================================================

[TO BE WRITTEN — discusses FPGAs, photonic integrated circuits, and
the practical implication of the EML Sheffer: a ``single-instruction
math co-processor'' that evaluates only $\eml(a, b)$ in hardware. The
point is that the formal proof of completeness implies that such a
co-processor is universal for the $\F_{36}$ class, modulo the boundary
points; this is a substantially stronger statement than the empirical
verification provides.]

% =====================================================================
\section{Conclusions and outlook}
\label{sec:conclusions}
% =====================================================================

[TO BE WRITTEN — summarises the artefact's contribution, the
philosophical significance of the witness-family form for Path~$C'$,
the open status of the seven SI questions, and the call-to-action
for the broader physics community: the existence of continuous
Sheffer operators is a structural fact about elementary functions
that deserves wider notice.]

\begin{thebibliography}{9}

\bibitem{odrzywolek2026eml}
A.~Odrzywołek,
\emph{All elementary functions from a single binary operator},
arXiv:2603.21852, 2026.
\url{https://arxiv.org/abs/2603.21852}

\bibitem{leanprover2021}
L.~de~Moura, S.~Ullrich,
\emph{The Lean 4 Theorem Prover and Programming Language},
in: Proceedings of CADE 2021, LNCS 12699, Springer, 2021, pp.~625--635.

\bibitem{harmonic2025aristotle}
Harmonic AI,
\emph{Aristotle: an automated mathematical reasoning service},
\url{https://harmonic.fun/aristotle}, accessed 2026.

\bibitem{naskrecki2026notes}
B.~Naskręcki,
\emph{Structure of the EML completeness proof: an expository tour},
companion notes \texttt{notes/proof\_structure.pdf} in the
\href{https://github.com/nasqret/eml-formalization}{eml-formalization}
repository, 2026.

\end{thebibliography}

\end{document}
